\emph{The space and its coordinates.} At $m>d$ one has $\varepsilon(m)\leq c/2<1$, so
\eqref{eq:doubling_cut} divides into $\lambda_m=\sum_{j\in W(m)}q_m(j)\lambda_j$ with $q_m>0$
and every index of $W(m)$ smaller than $m$. That identity is linear, so the sequences
satisfying it at every $m>d$ form a real vector space, and by induction it determines the whole
sequence from $\lambda_1,\dots,\lambda_d$: restriction to the first $d$ coordinates is a
linear isomorphism onto $\mathbb R^{d}$. The same induction, $q_m$ being positive, identifies
the positive sequences with the open positive orthant.

\emph{Linearity of $C$.} Each $\lambda\mapsto\lambda_m\,m^{p_*}$ is a linear functional of
$(\lambda_1,\dots,\lambda_d)$, and on the positive cone
Theorem~\ref{theo:doubling_sharp} makes their limit $C$ exist; additivity and positive
homogeneity therefore pass to $C$ on that cone. The cone being open and non-empty, every
element of the space is a difference of two of its elements, so $C$ extends unambiguously to a
linear functional on the whole space, and \eqref{eq:doubling_constant} is that functional
written in the coordinates $\lambda_1,\dots,\lambda_d$.

\emph{The signs.} At $\lambda_1=\dots=\lambda_d=1$ the identity reads $\sum_j\nu_j=C$, which is
positive because Theorem~\ref{theo:doubling_decay} confines $C$ to $[c_1,c_2]$ with $c_1>0$.
Fixing $j$ and letting $\lambda_j=t$ with the other coordinates $1$ makes $C$ an affine
function of $t$, positive at every $t>0$ by the same two theorems; an affine function of
negative slope is eventually negative, so $\nu_j\geq0$. Finally if $2j\leq d$ then every index
of every window $W(m)$ with $m>d$ exceeds $d/2\geq j$, so the induction never reads
$\lambda_j$: neither $\lambda_m$ for $m>d$ nor $C$ depends on it, and $\nu_j=0$.
